% ===== §2 The Legacy and Limits of Existentialism =====
\section{The Legacy and Limits of Existentialism}
\label{p14:sec:existentialism}

\noindent
The existentialist tradition is, among the great movements of modern philosophy, the one most seriously committed to the question of what it means to exist as a concrete, situated, finite being---as a being thrown into a world it did not choose, toward a death it cannot escape, in a body that is not merely its instrument but its very form of being-in-the-world. For this reason, it is the tradition from which any philosophy of shared happiness must begin. And for the same reason, it is the tradition whose limits are most instructive: because it is precisely where existentialism is most powerful---in its account of individual finitude and contingency---that it is least equipped to say what the flower on the path, shared with the one you love, actually is.

\subsection{Heidegger: Thrownness, the Moment of Vision, and the Limits of \emph{Mitsein}}
\label{p14:subsec:heidegger}

\noindent
Heidegger's fundamental ontology begins with a gesture that is still, after nearly a century, philosophically radical: the refusal to start with a subject who then encounters a world \parencite{heidegger1962}. \emph{Dasein}---the being for whom being is an issue---is always already in the world, always already thrown (\emph{geworfen}) into a situation not of its choosing, always already ahead of itself in its care (\emph{Sorge}) for its own being. The world is not an external reality that consciousness represents; it is the always-already-there within which \emph{Dasein} finds itself.

Within this framework, contingency receives a rigorous treatment. The flower on the path is, in Heideggerian terms, a piece of \emph{Zeug}---equipment, ready-to-hand (\emph{zuhanden}) in the first instance, revealing its worldly character only when it breaks down or solicits attention in a way that interrupts the flow of absorbed coping. When the flower arrests the walker's movement, it withdraws from the ready-to-hand and becomes present-at-hand (\emph{vorhanden}), an object of contemplation. More deeply, the flower participates in the structure of \emph{Befindlichkeit}---attunement or mood---which is the way in which \emph{Dasein} always already finds itself disposed toward the world before any explicit cognition. The walker who notices the flower does so in a mood, and the flower's beauty is not a property of the flower alone but the joint product of the flower and the attunement through which it is encountered.

So far, so illuminating. But what of the shared encounter? Heidegger does, of course, have a concept of being-with: \emph{Mitsein}, being-with-others, is part of the existential structure of \emph{Dasein} \parencite{heidegger1962}. \emph{Dasein} is not first a solitary subject who then enters into relations with others; it is constitutively with-others, and the others are not objects it encounters but fellow \emph{Daseins} whose being-in-the-world is always already co-constitutive of its own. The concept is important, and it marks a genuine advance over any simple subject-object model of intersubjectivity.

And yet \emph{Mitsein} is not enough for what the shared flower requires. The difficulty is structural: in Heidegger's account, \emph{Mitsein} is an existential---a necessary feature of \emph{Dasein}'s ontological structure---but it remains \emph{Dasein}'s structure. The others are co-constitutive of \emph{my} being-in-the-world; they enter into \emph{my} care-structure; they shape \emph{my} attunement. The fundamental unit of analysis remains \emph{Dasein}---and \emph{Dasein} is always mine (\emph{je meines}): ``the being of which this entity is an issue is in each case mine'' \parencite{heidegger1962}. \emph{Mitsein} is a structural feature of a being that remains fundamentally individual; it does not displace the individual as the locus of experience and event.

The moment of vision (\emph{Augenblick})---that authentic temporality in which \emph{Dasein} grasps its situation in the fullness of its thrownness and projection---is likewise \emph{Dasein}'s moment. It is the moment in which \emph{I} achieve authentic self-understanding, in which \emph{my} ownmost possibilities are disclosed. The concept has the precision and intensity that the shared flower seems to call for---that sudden illumination, that sense of the present moment held in its fullness---but it remains, in Heidegger's framework, an achievement of the individual \emph{Dasein}, not of the relational field between two \emph{Daseins}.

\begin{claim}[The Heideggerian limit]
Heidegger's \emph{Mitsein} establishes that being-with-others is ontologically constitutive of \emph{Dasein}. But it does not establish that there is an ontological unit---a \emph{we-ness}---that is prior to or irreducible to the individual \emph{Daseins} who compose it. The relational field remains, in Heidegger, a feature of \emph{Dasein}'s structure rather than a structure in its own right.
\end{claim}

\subsection{Sartre: Radical Contingency and the Hell of the Other}
\label{p14:subsec:sartre}

\noindent
If Heidegger's account of contingency is embedded in the structure of thrownness, Sartre's is more radical still \parencite{sartre1943}. For Sartre, existence precedes essence: there is no pregiven nature or telos that determines what a human being is or ought to be. The human being is condemned to be free---condemned because it did not choose this freedom, and yet cannot escape it; every situation, however constraining, leaves a residue of choice, and that residue is absolute. Contingency, for Sartre, goes all the way down: not just the situation into which one is thrown, but the very fact of one's existence is radically contingent, \emph{de trop}---superfluous, unjustified, in excess of any reason.

This is a profound and in many ways accurate description of something. The flower on the path is, in Sartrean terms, \emph{en-soi}---being-in-itself, brute, self-identical, without the self-transcendence that characterises human consciousness. The encounter with the flower is an encounter of \emph{pour-soi}---being-for-itself, the human consciousness that is always already beyond itself in its projects---with the in-itself. The beauty of the flower is not in the flower; it is the \emph{pour-soi}'s nihilating projection onto the opacity of being-in-itself.

But Sartre's account of the other---and here is what concerns us---is famously dark. The other is the gaze that turns me into an object: ``hell is other people'' (\zh{L'enfer, c'est les autres}) is not a casual remark but an ontological diagnosis \parencite{sartre1943}. The other threatens my freedom by constituting me as a facticity, a thing with determinate properties seen from the outside. The fundamental relation to the other is conflict: either I objectify the other or the other objectifies me, and love---Sartre is unsparing here---is the attempt to possess the other's freedom while keeping it free, which is a contradiction, and which therefore always fails.

This is not a framework in which shared happiness can find a home. The shared flower, in Sartre's account, is not an event in a relational field; it is a moment in which two freedoms, each radical and each threatening to the other, happen to direct their nihilating projections in the same direction. Whatever happiness arises is mine---my consciousness, my project, my freedom, temporarily and unstably aligned with another's. The sharing does not constitute a third thing; it is at best a momentary convergence of two absolute subjectivities that remain, at bottom, in competition.

\begin{claim}[The Sartrean limit]
Sartre's radical contingency and his account of freedom are philosophically indispensable. But his ontology of conflict between freedoms makes genuine relational co-constitution---the idea that the sharing of a contingent event might produce something irreducible to either participant---structurally impossible. The \emph{pour-soi} cannot share a world; it can only encounter other \emph{pour-sois} in a field of mutual threat and objectification.
\end{claim}

\subsection{Merleau-Ponty: Embodiment and the Flesh of the World}
\label{p14:subsec:merleau}

\noindent
Merleau-Ponty represents a significant departure from both Heidegger and Sartre, and one that moves---without fully arriving---in the direction this paper needs to go \parencite{merleau-ponty1962}. For Merleau-Ponty, the body is not the object of a constituting consciousness but the very medium of our being-in-the-world: perception is not representation but bodily engagement, and the perceived world is not a collection of objects but a field of solicitations and affordances structured by the body's capacities.

What this means for the encounter with the flower is significant. The beauty of the flower is not the \emph{pour-soi}'s nihilating projection onto an opaque in-itself; it is the response of a body-subject that is already attuned to the visible world through its motor habits, its perceptual history, its carnal engagement with things. The flower solicits the body; the body responds; and the beauty is in neither alone but in the encounter between a body with certain capacities and a world that calls forth those capacities.

More importantly for the shared encounter, Merleau-Ponty develops---especially in his later work on the ``flesh of the world'' (\emph{la chair du monde})---an account of intercorporeality: the idea that bodies are not isolated monads but participants in a shared carnal field, reversible touching-and-touched, seeing-and-seen \parencite{merleau-ponty1962}. When two embodied subjects perceive the same flower, they do so through bodies that are already in a kind of pre-personal intercorporeal relation, and the shared perception is not simply two separate perceptions happening simultaneously but a co-perception structured by this intercorporeal field.

This is genuinely closer to what the shared flower requires. The embodied resonance that neuroscience will later document---the mirror neuron system's cross-individual activation, the physiological synchrony of co-present bodies---has its philosophical anticipation in Merleau-Ponty's intercorporeality. And yet even here the relational field does not fully achieve ontological priority: the flesh of the world is a field of perception, and the co-perception of the flower is not yet the relational co-evolution that the GRB framework will identify as the locus of happiness. Intercorporeality is a structure of shared embodiment; it is not yet an account of how a contingent event descends upon a relational system and triggers the co-evolutionary dynamics in which happiness consists.

\subsection{Buber: The \emph{I--Thou} and the Threshold Not Crossed}
\label{p14:subsec:buber}

\noindent
Of all the existentialist and phenomenological thinkers, it is Martin Buber who comes closest to what this paper requires---and whose proximity to the goal makes his ultimate limitation all the more instructive \parencite{buber1970}.

Buber's fundamental distinction between the \emph{I--It} and the \emph{I--Thou} attitudes is well known. In the \emph{I--It} attitude, the other is an object in my world, to be used, analysed, represented; the relation is asymmetric, the other an \emph{It} that I constitute as an object of experience. In the \emph{I--Thou} attitude, the other is not an object but a subject who addresses me and to whom I respond; the relation is genuine meeting (\emph{Begegnung}), in which both parties are transformed and neither can be reduced to what they were before the meeting. ``In the beginning is the relation'' (\emph{Im Anfang ist die Beziehung})---this is Buber's most radical claim, and it is the one that points most directly toward the GRB framework: the relation is not secondary to the relata but constitutive of them.

And yet Buber does not complete the move that GRB requires. The \emph{I--Thou} relation, for Buber, is a relation \emph{between} an I and a Thou: it constitutes both, but both remain in some sense prior to the relation that constitutes them. The I that enters the \emph{I--Thou} relation is still, in Buber's account, an I that \emph{has} this relation---that can also fall back into the \emph{I--It} attitude, that can address or fail to address the other as Thou. The relation is the highest form of human existence for Buber, but it is not the \emph{ontological ground} of the subject; it is rather the subject's highest possibility.

Furthermore, Buber's account of the \emph{I--Thou} is primarily vertical---the meeting of the human I with the eternal Thou, God, as the ground of all genuine meeting---and while the horizontal meeting of two human beings is central to his thought, the relation between two finite \emph{I}s is always mediated by and pointing toward the eternal Thou. This gives the \emph{I--Thou} meeting a transcendent orientation that the GRB framework, with its immanent account of relational dynamics, does not share.

\begin{claim}[The Buberian threshold]
Buber's ``in the beginning is the relation'' is the closest approach, within the existentialist and phenomenological tradition, to the ontological claim that the GRB framework requires. But Buber does not complete the move: the relation remains, for him, the highest possibility of subjects who are in some sense already there to enter it, rather than the generative ground from which subjects emerge. The threshold that GRB crosses---locating the subject's very constitution in the relational field---Buber approaches but does not step over.
\end{claim}

\subsection{The Shared Diagnosis}
\label{p14:subsec:existdiagnosis}

\noindent
Across these four figures---and one could extend the survey to Jaspers's \emph{Existenzkommunikation}, to Levinas's ethics of the face, to the phenomenology of intersubjectivity in Husserl and Schütz---the same structural feature appears: the existentialist tradition, in all its richness, takes the individual subject as its fundamental unit, and adds the other, the relation, the co-presence as a modification of, or supplement to, or structure within, that individual subject's existence. Even Buber, who comes closest to reversing this priority, does not fully achieve the reversal.

The consequence for the philosophy of happiness is decisive. If the subject is always already individual, then happiness---whatever its content---is something that happens \emph{in} the subject: in its attunement, its freedom, its bodily engagement, its meeting with the Thou. The other, the relation, the sharing, can enrich or constitute the conditions for this happiness; but the happiness itself remains, in the last analysis, \emph{mine}. The shared flower, on this account, is either two separate happinesses occurring simultaneously, or one happiness (mine) that happens to require the other's presence as its condition. What it cannot be---within the existentialist framework---is a happiness that is \emph{irreducibly} of the relational field, that cannot be located in either participant because it does not reside in either.

This is precisely what the GRB framework will claim. But to understand what that claim requires, we must first examine the other great tradition this paper must engage and move beyond: the eudaimonist tradition, with its rich account of what flourishing consists in---and its equally systematic assumption that flourishing is the flourishing of an individual.
